%%=============================================================================
%% CAPITULO 1 -- INTRODUCAO
%%
%% Na introducao, faca a apresentacao do trabalho. Ela devera ocupar uma ou,
%% no maximo, duas paginas. De um panorama da tematica, exponha a questao
%% problema que conduzira a pesquisa proposta, a hipotese e a justificativa.
%%
%% Titulos: nao utilize ponto apos a numeracao das secoes.
%%=============================================================================

\section{Introdução}

Na primeira página da parte textual do trabalho de conclusão de curso, faça a
apresentação de sua pesquisa. A introdução deverá ocupar uma ou, no máximo, duas
páginas do seu trabalho. Nela, dê um panorama da temática, exponha a questão
problema que conduzirá a pesquisa proposta, a hipótese e a justificativa.

Os capítulos (seções) e subcapítulos (subseções) são sequenciais. O título de
cada seção primária (\textbf{1 INTRODUÇÃO}, \textbf{2 REFERENCIAL TEÓRICO} etc.)
é grafado em letras maiúsculas e negritadas. Não utilize ponto após a numeração
das seções. Deixe uma linha em branco entre as seções.

O Sistema Nacional de Pesquisa de Custos e Índices da Construção Civil (SINAPI) é a principal referência nacional de custos para obras e serviços de engenharia. Sua estrutura analítica e territorial, atualizada mensalmente, torna-o uma base pública singular para estudos de ciência de dados aplicados à política pública.
Em junho de 2026 a CAIXA lançou o desafio institucional de modernização do SINAPI por meio do certame 0071/2026, que aponta para a necessidade de selecionar empresas e startups de tecnologia capazes de desenvolver soluções modernas para ampliar, integrar e facilitar o acesso e uso dos dados da construção civil no Brasil. Essa diretriz converge diretamente com a proposta desta pesquisa, que busca aplicar modelos estatísticos espaço-temporais, machine learning e big data para apoiar decisões em políticas públicas de infraestrutura.
Apesar de amplamente utilizado, o SINAPI permanece predominantemente consultivo, sem exploração plena de seu potencial analítico para previsão e monitoramento. A pergunta central é: Porque é importante prever os dados do SINAPI?
Essa riqueza, entretanto, convive com diversos limites metodológicos como, a base pública não evidencia variáveis de execução física de obras, prazo de contrato, qualidade final do empreendimento ou desempenho físico-financeiro consolidado e o uso dos valores da região de São Paulo para as demais regiões quando da ausência de valores.
Por fim, essa obra tem como objetivo geral desenvolver e validar modelos de inteligência artificial e ciência de dados aplicados ao SINAPI para apoiar a tomada de decisão em políticas públicas de infraestrutura, como objetivos específicos responder as perguntas porque é importante prever os dados do SINAPI, e; na ausência dos dados do SINAPI em outras regiões, ao realizar a sobreposição pelos dados de São Paulo pode gerar um sobrepreço nas contratações públicas de construção civil?


\subsection{Objetivo geral}

Desenvolver e validar modelos de inteligência artificial e ciência de dados aplicados ao SINAPI para apoiar a tomada de decisão em políticas públicas de infraestrutura, possibilitando a previsão de custos da construção civil e a avaliação de distorções regionais que possam impactar a eficiência e a economicidade das contratações públicas.

\subsection{Objetivos específicos}

\begin{itemize}
  \item Objetivo específico 1 Analisar a importância da previsão dos dados do SINAPI para o planejamento orçamentário, a gestão de riscos e a formulação de políticas públicas de infraestrutura, identificando os benefícios da utilização de técnicas preditivas na tomada de decisão governamental;
  \item Objetivo específico 2 Desenvolver e avaliar modelos estatísticos espaço-temporais, de machine learning e de inteligência artificial aplicados aos dados do SINAPI, visando estimar a evolução futura dos custos da construção civil e identificar padrões regionais e temporais relevantes;
  \item Objetivo específico 3 Investigar os impactos da utilização de valores de referência do SINAPI da região de São Paulo para localidades sem dados próprios, avaliando o potencial de geração de sobrepreço e seus efeitos sobre a economicidade das contratações públicas de obras e serviços de engenharia.
\end{itemize}

\subsection{Problema}

Como modelos de inteligência artificial, machine learning e análise estatística espaço-temporal aplicados aos dados do SINAPI podem apoiar a tomada de decisão em políticas públicas de infraestrutura, por meio da previsão de custos da construção civil e da identificação de possíveis distorções decorrentes da utilização de referências de preços entre diferentes regiões do Brasil?

\subsection{Hipótese}

Parte-se da hipótese de que a aplicação de modelos estatísticos espaço-temporais, técnicas de machine learning e métodos de inteligência artificial aos dados do Sistema Nacional de Pesquisa de Custos e Índices da Construção Civil (SINAPI) é capaz de produzir previsões confiáveis da evolução dos custos da construção civil, contribuindo para o aprimoramento do planejamento orçamentário, da gestão de riscos e da formulação de políticas públicas de infraestrutura. Assume-se ainda que a utilização de valores de referência da região de São Paulo para localidades sem dados próprios pode introduzir distorções nos orçamentos de obras públicas, resultando em potenciais sobrepreços e reduzindo a eficiência na alocação dos recursos públicos.

\subsection{Justificativa}

A presente pesquisa justifica-se pela crescente necessidade de utilização de técnicas avançadas de ciência de dados e inteligência artificial para aprimorar a formulação e a implementação de políticas públicas baseadas em evidências. O Sistema Nacional de Pesquisa de Custos e Índices da Construção Civil (SINAPI) constitui a principal referência nacional para orçamentos, fiscalização e acompanhamento de obras públicas, sendo amplamente utilizado por órgãos governamentais, instituições financeiras, empresas e profissionais do setor da construção civil. Apesar de sua relevância estratégica, o uso do SINAPI ainda se concentra predominantemente na consulta de dados históricos, havendo espaço para a exploração de seu potencial analítico e preditivo.

Do ponto de vista acadêmico, a pesquisa contribui para a integração entre as áreas de ciência de dados, aprendizado de máquina, inteligência artificial, estatística espaço-temporal e administração pública, propondo a aplicação de métodos contemporâneos de análise de dados em uma base pública de grande abrangência nacional. Além disso, busca ampliar a literatura sobre o uso de técnicas preditivas na gestão pública, especialmente no contexto de infraestrutura e investimentos governamentais.

Sob a perspectiva social e institucional, a pesquisa possui potencial para contribuir com a melhoria da alocação de recursos públicos, o planejamento orçamentário, a antecipação de riscos em contratos de obras e a avaliação de distorções regionais de preços da construção civil. A capacidade de prever tendências de custos pode auxiliar gestores públicos na tomada de decisões mais eficientes, reduzindo desperdícios e aumentando a efetividade das políticas públicas de infraestrutura, habitação e saneamento.

A escolha desse recorte tem motivação prática e científica. Prática, em razão da experiência institucional relacionada ao SINAPI e da relevância crescente do tema para a modernização dos processos decisórios no setor público. Científica, pela oportunidade de investigar se modelos de inteligência artificial e machine learning podem gerar previsões confiáveis e produzir evidências sobre possíveis impactos da utilização de referências de preços de uma região em localidades com limitada disponibilidade de dados, contribuindo para o aprimoramento da gestão pública baseada em dados.
